\chapter{Intracranial Pressure Monitoring (Ventriculostomy)}

\section*{Overview}
Intracranial pressure (ICP) monitoring is the continuous measurement of pressure within the cranial vault using a catheter placed in the lateral ventricle or a sensor placed in the brain parenchyma or subdural space. An intraventricular catheter (ventriculostomy) also permits cerebrospinal fluid (CSF) drainage to reduce ICP.

\section*{Alternative Names}
ICP monitoring; intraventricular catheter; external ventricular drain; EVD; ventriculostomy; ICP bolt

\section*{Principle}
The intracranial contents (brain, blood, and CSF) are contained within the rigid skull, so pressure rises as volume increases. A catheter or sensor transmits pressure to an external transducer, and an intraventricular drain additionally allows CSF to be diverted to lower the pressure.

\section*{History}
Continuous ICP monitoring was pioneered by Nils Lundberg in the 1960s, and external ventricular drains became standard for monitoring and treating raised ICP in severe traumatic brain injury and other conditions.

\section*{Why It Is Done}
Ordered in patients with severe traumatic brain injury (typically Glasgow Coma Scale score of 8 or less with abnormal imaging), hydrocephalus, intracranial hemorrhage, cerebral edema, or after neurosurgery to guide management and maintain cerebral perfusion pressure.

\section*{Is This a Primary or Secondary Test or Procedure}
A primary monitoring and therapeutic procedure in neurocritical care for patients with actual or suspected raised intracranial pressure.

\section*{How It Is Performed}
Under sterile technique, usually at the bedside with or without a burr hole, a catheter is inserted into the frontal horn of the lateral ventricle or a sensor into the brain parenchyma. The system is connected to a pressure transducer zeroed at the level of the external auditory canal, and a CSF collection chamber is set at the desired drainage height.

\section*{Preparation}
Informed consent from the patient or surrogate, review of the coagulation profile with correction of coagulopathy, and confirmation that imaging shows no contraindication to the planned insertion site are required.

\section*{Contraindications and Precautions}
Coagulopathy, sepsis, and scalp infection at the insertion site are contraindications, while a small or shifted ventricle makes placement more difficult. Precautions include strict sterile technique, confirming catheter position with imaging, and placing the drain so that overdrainage does not occur.

\section*{Risks}
Risks include intracranial hemorrhage along the catheter tract, catheter-related ventriculitis or meningitis, infection, overdrainage or obstruction of the catheter, and neurological injury.

\section*{Aftercare and Recovery}
ICP and cerebral perfusion pressure are monitored continuously, and CSF may be drained against a pressure threshold. Sedation, hyperosmolar therapy, or surgery may be added when pressures remain elevated.

\section*{Outcome and Follow-up}
Normal adult ICP is below 10 to 15 mmHg, and sustained elevation above 20 mmHg is generally treated. Low pressure may indicate overdrainage or a leak, while failure of treatment to control ICP is a poor prognostic sign.

\section*{Expected Results}
Intracranial pressure below 10 to 15 mmHg with a normal waveform in an adult, or age-appropriate values in children, indicates satisfactory intracranial compliance.

\section*{Follow-up}
The drain is weaned or removed once ICP is controlled, and serial imaging and neurological examination assess for hydrocephalus, hemorrhage, or the need for a permanent shunt. CSF samples may be sent for culture if infection is suspected.

\section*{When to Seek Medical Care}
Follow the procedure team's specific instructions. Seek urgent medical assessment for severe or worsening pain, uncontrolled bleeding, fever or signs of infection, trouble breathing, chest pain, fainting, new weakness or confusion, inability to urinate, or any rapidly worsening symptom. The appropriate warning signs depend on the procedure and the patient's medical condition.

\section*{References}
\begin{enumerate}
  \item MedlinePlus. Intracranial pressure (ICP) monitoring. U.S. National Library of Medicine; 2025.
  \item Current Medical Diagnosis and Treatment. McGraw-Hill; 2025.
\end{enumerate}

\clearpage
